\section{Results and Analysis} \label{section:results_and_analysis}

This section presents the results of the DFT vs MLP interface ranking comparison across 1,381 interfaces evaluated during Iteration 3. The goal is to determine whether MACE, a large machine-learned interatomic potential, preserves favourability rankings obtained via DFT across a wide range of material systems and interface types.

\subsection{Overview of Rank Agreement}

\subsubsection{Summary trends in DFT vs. MLP agreement}

Ranking consistency between DFT and MLP was assessed using Spearman's $\rho$ and Kendall's $\tau$. Overall correlation was strong: $\rho$ ranged from 0.75 to 0.95 and $\tau$ from 0.70 to 0.90 across the majority of interface types. Metrics were extracted from \texttt{stats.csv}, which aggregates statistics across the full dataset. Perfect rank preservation (i.e. $\rho = 1$) was observed in a subset of alloy systems, while degraded consistency was evident in Silicon interfaces.

\subsubsection{Visual and tabular representation of metric distributions}

A scatter plot comparing DFT and MLP rankings across all interfaces is shown in Figure~\ref{fig:rank-scatter}, with supporting metric distributions detailed in Table~\ref{tab:global-metrics}. Top-N overlap scores (e.g. Top-10) were 5 for DFT@DFT vs MLP@MLP and 8 for DFT@DFT vs DFT@MLP, indicating modest agreement in identifying favourable structures.

\subsection{Cross-Method Evaluation}

\subsubsection{Comparison of MLP-relaxed vs. DFT-evaluated structures}

Due to computational constraints, MLP-relaxed structures evaluated with DFT (MLP@DFT) were omitted from this study. The cross-method comparison therefore focuses on DFT-relaxed structures evaluated with MLP (DFT@MLP). These exhibited higher agreement with DFT@DFT than MLP@MLP, indicating that relaxation geometry plays a dominant role in ranking discrepancies. This suggests a degree of structure-energy decoupling.

\subsection{System-Specific Observations}

\subsubsection{Carbon Systems (C|X)}

\paragraph{Note on incomplete data due to ISCA job failures}

Carbon-rich systems were underrepresented due to DFT job failures and high compute cost. As a result, only preliminary statistics could be extracted from \texttt{stats\_lower\_C.csv}, with limited rank correlation insight. These data points are not sufficient to draw conclusions but are retained for completeness.

\subsubsection{Silicon Interfaces}

\paragraph{Poor MLP rank preservation}

Silicon interfaces showed the weakest performance overall. Rank correlations were lower than for other systems, with $\rho < 0.70$ and $\tau < 0.65$ in many Si|Si and Si|X comparisons (see \texttt{stats\_lower\_Si.csv}). Errors tended to cluster in the lower half of the ranking, indicating failures in correctly identifying unfavourable interfaces.

\paragraph{Potential reason}

This poor agreement may stem from limited model generalisability across strained or undercoordinated Si environments. The MACE model was not explicitly fine-tuned for such systems, suggesting underfitting or domain mismatch.

\subsubsection{Germanium and Tin Interfaces}

\paragraph{Stronger correlation and better Top-N recovery}

Ge and Sn-containing systems yielded improved correlation metrics and Top-N recovery compared to Si. While Tin interfaces (\texttt{stats\_lower\_Sn.csv}) showed larger absolute energy errors, the rank preservation was generally robust. No catastrophic rank inversions were observed.

\subsubsection{Phase boundaries}

\paragraph{Overall poorer rank preservation}

Interfaces near the transition from favourable to unfavourable status (i.e. phase boundaries) were more susceptible to ranking shifts. While many transitions were preserved, files such as \texttt{stats\_lower\_Si\_upper\_Sn.csv} and \texttt{stats\_lower\_Ge\_upper\_Sn.csv} show increased rank distance near critical cutoffs.

\paragraph{Potential Reason}

These marginal cases are sensitive to small energy differences, which may be within the error tolerance of the MLP. This underscores the challenge of using MLPs for binary decision-making near phase boundaries.

\subsubsection{Perfect Alloys}

\paragraph{Perfect rank preservation}

Si|Ge alloy systems exhibited near-perfect rank alignment between DFT and MLP evaluations. Both $\rho$ and $\tau$ approached 1.0 across the board, as seen in \texttt{stats\_perfect\_alloys.csv}.

\paragraph{Potential Reason}

These systems feature uniform chemical environments with little interfacial disruption, meaning the atomic configurations fall well within the trained domain of the MACE model. The result is highly stable rank preservation.

\subsection{DFT vs MLP: Computational Cost Analysis}

MLP evaluations offered orders-of-magnitude speedups over DFT while retaining acceptable accuracy in most ranking tasks. However, the quality of this speedup is system-dependent and less reliable for Silicon-rich or mixed-coordination interfaces. These observations motivate a need for improved domain coverage, possibly via delta-learning or retraining on interfacial datasets.